\section{Purity of Mind}

From ``Power as Mind'' in \textit{The World as Power} by \textbf{Sir John Woodroffe}:

\begin{quotationx}
Mind reveals Consciousness by degrees, some minds more than the rest. The purer the mind the more it reflects or manifests Consciousness. The object then of the self-realising discipline or Sadhana is to purify the Mind so that it may manifest Consciousness. Purity of Mind is therefore to be sought. ``Pure'' and ``Purity'' are not used in their sexual sense only. This is only one and an elementary form of purity. It is obvious that if a Mind is dominated by sensuous desires and images, it cannot reflect or show Spirit. For this reason the Tantras in specifying the qualifications of the proposed disciple exclude the lewd and the glutton. It must be pure also in respect of other matters, and therefore free of greed, anger, envy and all else which is the mark of the impure Mind. Such a Mind is incapable of understanding spiritual things. But the Mind must not only be pure in the sense of freedom from what is bad, but must be positively kind and good and free from error. Purification of Mind is called \emph{Citta-buddhi}. The Mind must be an efficient and trained instrument of knowledge which is its appropriate food, and should if necessary be sharpened by the study of logic and the practice of debate. It should be made capable in this and other ways of understanding the highest metaphysical ideas. And so the disciple is recommended to study the sacred texts, Logic and Metaphysic. At the same time there should be devotion to and worship of God as the Mother-Power (one with Siva as unchanging Consciousness) who is called Lalita, Mahakali Maha-tripurasundari, Mahakundalini and by other names which denote only aspects of the one Reality as Power.

Ritual is the art of religion. The rituals are designed to secure realisation of Unity with Her. \emph{Sakta-Sadhana} which term includes what is called in English ``ritual'', is based on sound psychological principles with which I will deal in another volume. The ritual is an expression in action of the philosophical principles above described. Thus the whole evolving cosmic process is imagined in the rite called \emph{Bhuta-suddhi}, in which each of the lower principles is merged in the higher, until in imagination the abode of \emph{Siva-Sakti} is reached. So also the \emph{Sri-Yantra} or Diagram represents both the body of the Sadhaka as the Microcosm and the whole universe. All ordinary acts and functions become worship by dedication to the Mother-Power, and self-identification with that Power in all physical functions and acts. The \emph{Sadhana} then realises himself as the Mother-Power in the form of himself.
\end{quotationx}



\flrightit{Posted on 2008-06-11 by Anonymous}

